\section{Quintic normal forms and exact Tschirnhaus certificates}\label{sec:normal-forms}

\subsection{Depressing the general quintic}

Consider
\begin{equation}\label{eq:general-quintic}
 A x^5+B x^4+C x^3+D x^2+E x+F=0,
 \qquad A\ne0.
\end{equation}
Set \(x=y-B/(5A)\) and divide by \(A\).  Direct expansion gives a depressed
quintic
\begin{equation}\label{eq:depressed-quintic}
 y^5+p y^3+q y^2+r y+s=0,
\end{equation}
where
\begin{align}
 p&=\frac{5AC-2B^2}{5A^2},\label{eq:depress-p}\\
 q&=\frac{25A^2D-15ABC+4B^3}{25A^3},\label{eq:depress-q}\\
 r&=\frac{125A^3E-50A^2BD+15AB^2C-3B^4}{125A^4},\label{eq:depress-r}\\
 s&=\frac{3125A^4F-625A^3BE+125A^2B^2D-25AB^3C+4B^5}
          {3125A^5}.\label{eq:depress-s}
\end{align}

\begin{theorem}[Exact depression identity]\label{thm:depress-quintic}
Equations~\eqref{eq:depress-p}--\eqref{eq:depress-s} satisfy the identity
\begin{equation}\label{eq:depress-identity}
 \frac1A P\!\left(y-\frac{B}{5A}\right)
 =y^5+py^3+qy^2+ry+s.
\end{equation}
After multiplication by \(3125A^5\), this is an integer polynomial identity
and hence a theorem of \ipra.
\end{theorem}

\begin{proof}
Expand the left side, collect powers of \(y\), and use the binomial
coefficients \(1,5,10,10,5,1\).  The coefficient of \(y^4\) is
\(-B/A+B/A=0\); the remaining coefficients normalize to the four displayed
fractions.  The executable verifier checks the fully symbolic numerator.
\end{proof}

The formal derivative of~\eqref{eq:general-quintic} is a quartic.  Thus all
critical points are roots of a polynomial whose factorization identity is
already covered by the checked Ferrari module.  For real rational coefficients,
a Sturm chain or derivative-based interval subdivision is a bounded exact
algorithm for isolating every real root.  This gives an independent effective
meaning to the source page's elementary observation that a quintic's turning
points are controlled by a quartic.

\subsection{A transparent reduction to principal form}

A principal quintic has the form
\begin{equation}\label{eq:principal-form}
 Y^5+c_2Y^2+c_1Y+c_0=0;
\end{equation}
its quartic and cubic terms vanish.  We now prove the quadratic Tschirnhaus
step by power sums.

Let \(x_1,\ldots,x_5\) be the roots of a monic depressed quintic in a formal
splitting algebra, and set
\[
 S_m=\sum_{i=1}^5x_i^m.
\]
Then \(S_1=0\).  If the coefficient \(p\) of \(x^3\) in
\eqref{eq:depressed-quintic} is zero, the equation is already principal.  Assume
therefore \(p\ne0\).  Newton's identity gives \(S_2=-2p\ne0\).

Define
\begin{equation}\label{eq:quadratic-tsch}
 Y_i=x_i^2+\alpha x_i+\beta,
 \qquad \beta=-\frac{S_2}{5}.
\end{equation}
Then
\begin{equation}\label{eq:Ysum}
 \sum_iY_i=S_2+5\beta=0.
\end{equation}
A second expansion gives
\begin{align}
 \sum_iY_i^2
 &=S_4+2\alpha S_3+(\alpha^2+2\beta)S_2+5\beta^2\notag\\
 &=S_4+2\alpha S_3+\alpha^2S_2-\frac{S_2^2}{5}.
 \label{eq:Ysquare-sum}
\end{align}
Choose \(\alpha\) to satisfy
\begin{equation}\label{eq:alpha-principal}
 5S_2\alpha^2+10S_3\alpha+5S_4-S_2^2=0.
\end{equation}
This is a genuine quadratic because \(S_2\ne0\), and its required square root
is supplied by the checked quadratic factorization theorem.

\begin{theorem}[Principal-form reduction]\label{thm:principal-reduction}
With \(\alpha,\beta\) satisfying
\eqref{eq:quadratic-tsch}--\eqref{eq:alpha-principal}, the transformed roots
\(Y_i\) are the roots of a polynomial of the form~\eqref{eq:principal-form}.
The transformation and its correctness are expressible by finite radical-tower
certificates and are provable in \ipra.
\end{theorem}

\begin{proof}
Equations~\eqref{eq:Ysum} and~\eqref{eq:Ysquare-sum} show that the first two
power sums \(T_1=\sum Y_i\) and \(T_2=\sum Y_i^2\) vanish.  If
\(e_j(Y)\) denotes the elementary symmetric polynomial in the \(Y_i\), Newton's
identities give
\[
 e_1=T_1=0,
 \qquad
 2e_2=e_1T_1-T_2=0.
\]
Characteristic zero permits cancellation of \(2\).  Therefore
\[
 \prod_{i=1}^5(Z-Y_i)
 =Z^5-e_3(Y)Z^2+e_4(Y)Z-e_5(Y),
\]
which is principal.  Every symmetric coefficient can either be computed as a
bounded sum of products or obtained as a resultant.  The proof uses only the
quadratic relation~\eqref{eq:alpha-principal} and polynomial identities.
\end{proof}

The two possible roots of~\eqref{eq:alpha-principal} give two valid principal
forms.  This explains constructively the source page's statement that either
solution may be used.  No choice principle is needed: a certificate carries one
chosen radical and its defining square equation.

\subsection{Bring--Jerrard reduction as a compact certificate}

Let
\[
 P(Y)=Y^5+c_2Y^2+c_1Y+c_0
\]
a principal quintic.  A quartic Tschirnhaus map has the form
\begin{equation}\label{eq:quartic-tsch}
 T(Y)=Y^4+\alpha Y^3+\beta Y^2+\gamma Y+\delta.
\end{equation}
Introduce a new indeterminate \(V\) and compute
\begin{equation}\label{eq:tsch-resultant}
 R(V)=\Res_Y\bigl(P(Y),V-T(Y)\bigr).
\end{equation}
Because the first polynomial has degree five, \(R\) is monic of degree five in
\(V\), counting multiplicities.  A successful Bring certificate consists of
\[
 (\alpha,\beta,\gamma,\delta,d_1,d_0).
\]
together with their finite algebraic-tower definitions and the exact identity
\begin{equation}\label{eq:bring-certificate}
 R(V)=V^5+d_1V+d_0.
\end{equation}

\begin{theorem}[Soundness of a Bring--Jerrard certificate]
\label{thm:bring-certificate-sound}
If the finite identity~\eqref{eq:bring-certificate} holds and \(y\) is a root
of \(P\), then \(v=T(y)\) satisfies
\[
 v^5+d_1v+d_0=0.
\]
Checking the certificate is primitive recursive, uniformly in the coded
coefficient tower, and its soundness is provable in \ipra.
\end{theorem}

\begin{proof}
The defining property of the resultant is already visible from the Sylvester
determinant: substituting \(V=T(y)\) makes the two polynomials
\(P(Y)\) and \(T(y)-T(Y)\) have the common root \(Y=y\), so the determinant
vanishes.  Substitution into~\eqref{eq:bring-certificate} gives the desired
trinomial equation.  Alternatively a Bézout certificate
\(A(Y,V)P(Y)+B(Y,V)(V-T(Y))=R(V)\) makes this a single polynomial evaluation.
All matrices and degree bounds are finite.
\end{proof}

\begin{frontier}[Universal witness generation]
The classical Bring construction chooses the four parameters in
\eqref{eq:quartic-tsch} so that the coefficients of \(V^4,V^3,V^2\) in
\eqref{eq:tsch-resultant} vanish.  After suitable elimination, the remaining
choices require a cubic and a quadratic equation.  This is entirely compatible
with \ipra{} because the Cardano and quadratic factorization identities are
already checked.  What is large is not the logical principle but the universal
symbolic witness: expanded formulas can occupy megabytes.  This packet proves
the uniform verifier and soundness theorem, and gives an exact route for adding
any generated witness, but it does not claim a kernel-checked universal
multimegabyte certificate generator.
\end{frontier}

This distinction is important.  The mathematically useful theorem is often the
small verifier~\eqref{eq:bring-certificate}, not a human-readable expansion of
all its witnesses.

\subsection{Scaling a Bring--Jerrard trinomial}

The one-variable reduction is short and complete.

\begin{theorem}[Scaling formula]\label{thm:bring-scaling}
Let \(c^4=p\) and \(c\ne0\).  Then
\begin{equation}\label{eq:scaling-identity}
 (cz)^5+p(cz)+q
 =c^5\left(z^5+z+\frac{q}{c^5}\right).
\end{equation}
Consequently every compatible branch of a root of \(x^5+px+q\) has the form
\begin{equation}\label{eq:root-by-BR}
 x=c\,\BR\!\left(\frac{q}{c^5}\right),
\end{equation}
with the other four roots obtained by the other algebraic branches.  In the
customary notation this is
\(p^{1/4}\BR(p^{-5/4}q)\).  If \(p=0\), the equation is simply
\(x^5=-q\).
\end{theorem}

\begin{proof}
Use \(p=c^4\) in the left side and factor out \(c^5\).  Formula
\eqref{eq:root-by-BR} follows from the defining equation of the selected Bring
branch.  To avoid division inside \ipra{}, multiply by \(c^5\) and retain the
guard \(c\ne0\).
\end{proof}

Choosing \(c^4=-p\) instead produces the alternative form
\(z^5-z+a=0\).  The four choices of the fourth root and the five algebraic
branches explain why the displayed scaling formula must be accompanied by
branch compatibility; it is not a single global complex-valued radical.

\subsection{Inverting the Tschirnhaus transformation}

Suppose a root \(v\) of the Bring--Jerrard equation has been computed.  The
corresponding original root \(y\) is a common zero of
\begin{equation}\label{eq:inverse-system}
 P(Y)=0,
 \qquad T(Y)-v=0.
\end{equation}
The polynomial gcd of these two polynomials over the finite coefficient tower
can be computed by the subresultant algorithm.  In the generic case the gcd
has degree at most four, so its roots are available from the checked formulas
through Ferrari.  In a degenerate case the certificate records the actual gcd
and its degree; no generic division is silently performed.

\begin{proposition}[Certified inversion]
A list of candidates \(y_j\) is a complete inverse image of \(v\) under the
Tschirnhaus map on the roots of \(P\) when the following finite conditions
hold:
\begin{enumerate}
\item every \(y_j\) satisfies both equations in~\eqref{eq:inverse-system};
\item the monic gcd of the two left sides factors as \(\prod_j(Y-y_j)\);
\item the supplied Bézout/subresultant trace verifies that this polynomial is
      the gcd.
\end{enumerate}
These conditions have a primitive-recursive verifier and imply completeness in
\ipra.
\end{proposition}

This is the constructive resolution of the ``extraneous solutions'' issue
mentioned in the source page: numerical guessing is not logically necessary.
One may approximate candidates to choose a branch, but exact substitution and
the gcd certificate perform the final verification.

\subsection{Brioschi and other normal forms}

The Brioschi family
\[
 W^5-10CW^3+45C^2W-C^2=0.
\]
afforded a different route in the classical literature.  The source page gives
a rational Tschirnhaus map
\[
 W=\frac{\lambda+\mu X}{X^2/C-3}.
\]
For any proposed formulas for \(C,\lambda,\mu\), correctness is again a finite
certificate: multiply by the fifth power of the denominator, reduce modulo the
original quintic, and check that the remainder is zero.  The denominator guard
is retained, and exceptional roots are checked separately.  Thus the Brioschi
normal-form claim has the same compact \ipra{} interface as
Theorem~\ref{thm:bring-certificate-sound}, even though deriving the universal
polyhedral formulas for the parameters belongs to the global icosahedral
frontier of Section~\ref{sec:beyond}.
